\documentclass[11pt,a4paper]{article}

\usepackage[T1]{fontenc}
\usepackage[utf8]{inputenc}
\usepackage{newpxtext,newpxmath}
\usepackage[protrusion=true]{microtype}
\usepackage[a4paper,left=2.2cm,right=2.2cm,top=1.8cm,bottom=1.8cm]{geometry}
\usepackage{xcolor}
\usepackage[hidelinks]{hyperref}

\definecolor{accent}{RGB}{28,58,104}
\definecolor{subtle}{RGB}{110,110,110}
\setlength{\parindent}{0pt}
\setlength{\parskip}{0.65em}
\linespread{1.04}
\pagestyle{empty}
\hypersetup{
  pdftitle={Cover Letter -- Postdoctoral Application},
  pdfauthor={Avi Kadria}
}

\begin{document}

{\Large\bfseries\color{accent}Avi Kadria}\hfill November 12, 2026\\[0.2em]
Department of Computer Science, Bar-Ilan University\\
{\small\color{subtle}\href{mailto:avi.kadria3@gmail.com}{avi.kadria3@gmail.com}
\quad\textbar\quad\href{https://avika81.github.io/papers/}{avika81.github.io/papers/}}

\vspace{0.6em}
\textbf{Application for a Postdoctoral Position in Theoretical Computer Science}

Dear Members of the Selection Committee,

I am applying for a postdoctoral position in theoretical computer science. I am a PhD candidate in Computer Science at Bar-Ilan University, where I study graph algorithms and fine-grained complexity. My research focuses on distance problems, including shortest paths, distance oracles, routing, girth, and diameter, and on the trade-offs between running time, space, and approximation.

My work combines the design of faster algorithms with lower bounds that identify the limits of existing methods. With Liam Roditty, I improved the construction time of Thorup--Zwick distance oracles (ICALP 2026). Our recent work on single-source shortest paths in weighted undirected graphs, accepted to SOFSEM 2027, gives the first running-time improvement since the FOCS 2023 result of Duan, Mao, Shu, and Yin. My work has also appeared at SODA, ESA, ISAAC, and DISC; our ISAAC 2025 paper received the Best Paper Award.

Recent projects extend these questions in several directions. I have developed a deterministic girth approximation algorithm, and my joint work studies lower bounds for pivot-based distance oracles in sparse graphs, randomized compact routing schemes, and new analytical tools for the Thorup--Zwick hierarchy. In ongoing work on weighted planar digraphs, we use geometric range searching to obtain a near-linear-time approximation scheme for diameter.

As a postdoc, I plan to investigate the complexity of shortest paths, improve the construction and space bounds of distance oracles, and understand when graph structure allows faster distance algorithms. I am particularly interested in developing tools that apply across these problems, and in using algorithmic progress to guide lower-bound questions. I aim to pursue these directions through collaborations with faculty, postdoctoral researchers, and students.

I received the Clore Scholarship in 2025 and the President Scholarship in 2024. Before returning to full-time research, I served as an officer in Unit 8200, where I established and led a cyber research and development team. These experiences shaped my approach to mentoring, organizing long-term projects, and working across different areas of expertise.

Thank you for considering my application. My CV, complete publication list, and research plan accompany this letter. I would welcome the opportunity to discuss my work and future research.

\vspace{0.4em}
Sincerely,\\[0.3em]
Avi Kadria

\end{document}
